\chapter{CONCLUSION AND RECOMMENDATION}
\section{Conclusions}
The comparative seismic assessment of the G+4 Infectious and Communicable Disease Hospital indicates that foundation flexibility has a significant influence on the response of its six structural blocks. From the fixed-base to the soft-soil condition, the fundamental period increases by 8.53\%, 19.73\%, 5.22\%, 8.99\%, 9.00\% and 7.85\% for Blocks 1--6, respectively. The first mode satisfying 90\% cumulative mass participation in both horizontal directions varies between Modes 5 and 10.

Uppermost-storey displacements in the X-direction response-spectrum cases increase as the foundation support becomes more flexible in all six blocks. The maximum reported inter-storey drift is 0.774\%, occurring in Block 3 under the soft-soil condition. The response-spectrum base shears, however, exhibit relatively small variations among the foundation-support conditions in several blocks. Hence, increasing the vibration period does not necessarily lead to a reduction in the seismic force demand.

Vertical foundation movements under the $1.2DL+1.5LL$ load combination generally become larger with increasing soil flexibility. The beam and column force envelopes also change with support stiffness, but the magnitude and direction of change vary with structural configuration and the member-force component. X-direction fragility curves supplement the assessment of performance under different soil conditions. Collectively, the results demonstrate the importance of incorporating soil--foundation flexibility in the seismic assessment of critical healthcare infrastructure.

\section{Recommendations}
The analysis should be extended through detailed verification of foundation dimensions and stiffness parameters for the actual site. Three-dimensional structural models should incorporate the appropriate soil properties, structural connections and design load combinations. Additional earthquake analyses and consistent fragility thresholds would strengthen assessment of damage probabilities in each structural block.

The observed differences between displacement, drift, shear and member-force response indicate that soil--structure interaction should be evaluated through several response measures rather than fundamental period alone. Particular attention should be given to lateral drift and foundation deformation in the more flexible-support cases.
